%\documentclass[journal]{IEEEtran}
\documentclass[journal,12pt,draftclsnofoot,onecolumn]{IEEEtran}
\IEEEoverridecommandlockouts

\usepackage{cite}
\usepackage{amsmath,amssymb,amsfonts,mathrsfs}
\usepackage{bbm}
\usepackage[normalem]{ulem}
\newtheorem{theorem}{Theorem}
\newtheorem{lemma}{Lemma}
\newtheorem{corollary}{Corollary}
\newtheorem{remark}{Remark}
\newtheorem{proposition}{Proposition}
\newtheorem{definition}{Definition}
\newtheorem{problem}{Problem}
\newtheorem{assumption}{Assumption}
\usepackage{placeins}
\usepackage{graphicx}
\usepackage{textcomp}
\usepackage{xcolor}
\def\BibTeX{{\rm B\kern-.05em{\sc i\kern-.025em b}\kern-.08em
    T\kern-.1667em\lower.7ex\hbox{E}\kern-.125emX}}


\begin{document}

\title{Finite blocklength bounds for composite binary hypothesis testing via R\'enyi divergences}
\author{\IEEEauthorblockN{El{\'i}as Vera-Sig{\"u}enza}\\
\IEEEauthorblockA{
\textit{Information Theory, Probability, and Statistics Unit}\\
\textit{Okinawa Institute of Science and Technology} \\
Okinawa, Japan\\
{email: elias.vera@oist.jp,} \\
ORCID: 0000-0003-4090-3557}}

\maketitle


\begin{abstract}
We derive finite blocklength converse and achievability bounds for minimax composite binary hypothesis testing under a uniform Type I error probability constraint. The converse reduces any admissible composite test to simple binary subproblems. For achievability, a joint R\'enyi projection induces a log-likelihood ratio that yields a single threshold test uniformly bounding the error probabilities. When evaluated under an exponentially decaying Type I error constraint, these finite bounds establish the infimum of the Kullback--Leibler divergence from the alternative class to the null class as the threshold rate separating exponential decay of the minimax Type II error probability from its exponential convergence to one. For finite alphabets, a saddle-point analysis yields the exact minimax Type II error exponent for all rates below this threshold, together with the matching polynomial order and logarithmic threshold correction. In the same finite-alphabet setting, under a fixed Type I error constraint, the limit of the rescaled R\'enyi divergence as its order approaches zero recovers the minimum Kullback--Leibler divergence from the null class to the alternative class as the exact minimax Type II error exponent. Finally, we establish sufficient conditions under which the projected pair is least favourable at finite blocklength, yielding an exact reduction of the composite problem to simple binary hypothesis testing.
\end{abstract}
\begin{IEEEkeywords}
composite binary hypothesis testing, minimax hypothesis testing, finite blocklength bounds, R\'enyi divergence, R\'enyi projection, Hoeffding exponent, strong converse, robust detection
\end{IEEEkeywords}

%##################################################
%##################################################
\section{Introduction}
\label{sec:introduction}

Binary hypothesis testing asks whether observed data were generated under a null hypothesis or an alternative hypothesis. In the Neyman--Pearson (NP) formulation, a test is required to control the Type I error; concretely, the rejection of the null hypothesis when it is true, whilst minimising the Type II error, or the acceptance of the null hypothesis when the alternative is true~\cite{neyman1933testing,lehmann2005testing}. 

When both hypotheses are simple, the NP lemma identifies an optimal likelihood ratio test at every blocklength $n$~\cite{neyman1933testing}. The associated asymptotic tradeoff is also classical. Under a fixed Type I error constraint, the Chernoff--Stein lemma identifies the Kullback--Leibler (KL) divergence as the optimal Type II error exponent~\cite{cover2006elements}. Moreover, if the Type I error is required to decay exponentially, the tradeoff exhibits a transition between a regime in which the Type II error decays exponentially and a strong converse regime in which it converges to one. The threshold rate separating these regimes is governed by the KL divergence from the alternative hypothesis to the null hypothesis~\cite{blahut1974hypothesis,han2002strong}.

Bruno et al.~\cite{bruno2026finite} recently derived finite blocklength R\'enyi achievability and converse bounds for the simple binary hypothesis testing problem, recovering this classical threshold rate and establishing exponential decay of the Type II error below it.

In this work, both hypotheses are composite. The null hypothesis and the alternative hypothesis each comprise a class of probability laws. A single test must satisfy the Type I error constraint uniformly over the null class, whilst its Type II error is evaluated through the supremum over the alternative class. The resulting problem is intrinsically minimax. A test designed for one fixed pair may satisfy the required Type I constraint and yield a small Type II error for that pair, whilst failing to provide either guarantee uniformly over the full classes.

Composite hypothesis testing has long been studied through robust detection and least favourable distributions. Under suitable conditions, a least favourable pair reduces the composite minimax problem exactly to a simple binary test~\cite{huber1965robust,verdu2003minimax,fauss2021minimax}. Such reductions are powerful, but they require additional structure. Recent work by G\"ul~\cite{gul2026structural} highlights the limitations of such reductions by constructing convex uncertainty classes for which a single pair governs the asymptotically optimal error exponent, yet fails to be least favourable at finite blocklength.

Other studies characterise error exponents for structured composite problems. Tomamichel and Hayashi~\cite{tomamichel2017operational} consider a fixed null distribution and structured composite alternatives. They derive the threshold rate, the optimal error exponent, the strong converse exponent, and second-order asymptotics. For product and Markov alternative classes, these exponents are expressed through Sibson-$\alpha$ mutual information and R\'enyi conditional mutual information. Shayevitz~\cite{shayevitz2011renyi} studies a composite problem with two sensors and relates the maximum miss-detection exponent to a R\'enyi divergence between two distribution families under an extremal noise condition. In the quantum setting, Berta, Brand\~ao, and Hirche~\cite{berta2021composite} study composite hypotheses formed from convex combinations of independent and identically distributed quantum states, characterising the maximum Type II error exponent by a regularised divergence. Exact asymptotic error exponents for classical composite hypotheses over convex classes are also known in broader formulations, including Mosonyi, Szil\'agyi, and Weiner~\cite{mosonyi2021error}.

\textit{These results characterise asymptotic performance}, but do not provide explicit finite blocklength minimax achievability and converse bounds for the classical composite versus composite problem considered here. In this light, we provide a finite blocklength treatment of the composite problem that does not rely on an exact reduction to a simple binary pair.


%##################################################
%##################################################
\section{Notation and problem formulation}
\label{sec:problem-formulation}

Let $\mathcal P$ and $\mathcal Q$ be nonempty classes of probability laws on a common measurable space $(\mathcal X,\mathcal F)$. For a blocklength $n\geq1$, let $X^n=(X_1,\ldots,X_n)$. The composite binary hypothesis testing problem is
\begin{equation*}
\begin{aligned}
H_{\mathcal P}
&:X^n\sim P^{\otimes n}
\quad\text{for some }P\in\mathcal P,\\
H_{\mathcal Q}
&:X^n\sim Q^{\otimes n}
\quad\text{for some }Q\in\mathcal Q.
\end{aligned}
\end{equation*}

A randomised test at blocklength $n$ is an $\mathcal F^{\otimes n}$--measurable function $\varphi_n:\mathcal X^n\to[0,1]$. For an observation $x^n$, the value $\varphi_n(x^n)$ is the probability of deciding $H_{\mathcal Q}$, whilst $1-\varphi_n(x^n)$ is the probability of deciding $H_{\mathcal P}$.

For a randomised test $\varphi_n$, the Type I and Type II error probabilities are
\begin{equation*}
\alpha_n(\varphi_n;\mathcal P):=\sup_{P\in\mathcal P}\mathbb E_{P^{\otimes n}}\left[\varphi_n(X^n)\right],
\end{equation*}
and
\begin{equation*}
\beta_n(\varphi_n;\mathcal Q):=\sup_{Q\in\mathcal Q}
\mathbb E_{Q^{\otimes n}}\left[1-\varphi_n(X^n)\right].
\end{equation*}
We consider the asymmetric setting in which the Type II error probability is minimised subject to a uniform constraint on the Type I error probability. For a specified threshold $\varepsilon\in(0,1)$, define
\begin{equation}
\beta_n^\star(\varepsilon;\mathcal P,\mathcal Q):=\inf_{\varphi_n:\,\alpha_n(\varphi_n;\mathcal P)\leq\varepsilon} \beta_n(\varphi_n;\mathcal Q),
\label{eq:composite-minimax}
\end{equation}
where the infimum is over randomised tests at blocklength $n$. Thus, $\beta_n^\star(\varepsilon;\mathcal P,\mathcal Q)$ is the minimum Type II error probability at blocklength $n$, evaluated over the alternative class, subject to the requirement that the Type I error probability does not exceed $\varepsilon$ uniformly over the null class. The threshold $\varepsilon$ may depend on the blocklength. In particular, the exponentially decaying Type I error constraint considered later is obtained by setting $\varepsilon=e^{-nr}$ for some $r>0$. For fixed probability laws $P$ and $Q$, we use the same notation for the corresponding simple binary problem:
\begin{equation*}
\beta_n^\star(\varepsilon;P,Q) := \inf_{\varphi_n:\, \mathbb E_{P^{\otimes n}}[\varphi_n(X^n)]\leq\varepsilon} \mathbb E_{Q^{\otimes n}}\left[1-\varphi_n(X^n)\right],
\end{equation*}
where the infimum is again over randomised tests at blocklength $n$. For $\lambda>0$, $\lambda\neq1$, let $\mu$ be any measure dominating $P$ and $Q$. The order--$\lambda$ Hellinger integral is
\begin{equation}
Z_\lambda(Q,P):=\int_{\mathcal X} \left(\frac{dQ}{d\mu}\right)^\lambda \left(\frac{dP}{d\mu}\right)^{1-\lambda} \,d\mu,
\label{eq:hellinger-integral}
\end{equation}
and the R\'enyi divergence of order $\lambda$ is
\begin{equation}
D_\lambda(Q\|P):=\frac{1}{\lambda-1}\log Z_\lambda(Q,P).
\label{eq:renyi-divergence-definition}
\end{equation}
These quantities do not depend on the choice of $\mu$. Standard extended-real conventions apply, and $D(Q\|P)$ denotes the KL divergence. For the composite classes, define
\begin{equation}
\begin{aligned}
D_\lambda(\mathcal Q\|\mathcal P) &:= \inf_{Q\in\mathcal Q}
\inf_{P\in\mathcal P} D_\lambda(Q\|P),\\
D(\mathcal Q\|\mathcal P) &:= \inf_{Q\in\mathcal Q} \inf_{P\in\mathcal P} D(Q\|P).
\label{eq:inf-KL}
\end{aligned}
\end{equation}
%###################################################
%###################################################
\section{Finite blocklength bounds}
\label{sec:finite-blocklength-bounds}

\subsection{Converse bound}
\label{sec:composite-converse}

We first derive a lower bound on $\beta_n^\star(\varepsilon;\mathcal P,\mathcal Q)$ using the R\'enyi divergence of order $\lambda>1$. 

%---------------------------------------------------
\begin{theorem}
\label{thm:finite-blocklength-composite-renyi-converse} For every $\varepsilon\in(0,1)$ and $n\geq1$,
\begin{equation}
\beta_n^\star(\varepsilon;\mathcal P,\mathcal Q)
\geq 1- \exp\left\{-\sup_{\lambda>1}
\frac{\lambda-1}{\lambda}\left[\log\frac{1}{\varepsilon}-nD_\lambda(\mathcal Q\|\mathcal P)
\right]_+\right\},
\label{eq:finite-blocklength-composite-renyi-converse}
\end{equation}
where $[a]_+:=\max\{a,0\}$ for $a\in\mathbb R$, with $[-\infty]_+:=0$. 
\end{theorem}
Proof of theorem~\ref{thm:finite-blocklength-composite-renyi-converse} is given in Appendix~\ref{app:finite-blocklength-composite-renyi-converse}.

%###################################################
%###################################################
\subsection{Achievability bound}
\label{sec:composite-achievability}

Achievability requires a single test that controls the Type I error uniformly over $\mathcal P$ and the Type II error uniformly over $\mathcal Q$. The following lemma gives sufficient conditions for a
threshold test to satisfy both requirements.

%---------------------------------------------------
\begin{lemma}
\label{lem:uniform-score-achievability} Fix $n\geq1$, $\varepsilon\in(0,1)$, $0<\lambda<1$, and $d\geq0$. Suppose that a measurable function $h:\mathcal X\to[-\infty,+\infty]$ satisfies
\begin{equation*}
\begin{aligned}
\sup_{P\in\mathcal P}
\mathbb E_P\left[e^{\lambda h}\right]
&\leq
e^{(\lambda-1)d},
\\
\sup_{Q\in\mathcal Q}
\mathbb E_Q\left[e^{(\lambda-1)h}\right]
&\leq
e^{(\lambda-1)d}.
\end{aligned}
\end{equation*}
For $\tau\in\mathbb R$, define the threshold test
\begin{equation*}
\psi_{n,\tau}(x^n):=\mathbbm 1\left\{\sum_{i=1}^n h(x_i)\geq\tau\right\},
\end{equation*}
where the sum is set to zero on sequences for which both $+\infty$ and $-\infty$ occur. Every threshold satisfying
\begin{equation*}
\tau \geq \frac{\log(1/\varepsilon)-n(1-\lambda)d}{\lambda}
\end{equation*}
gives
\begin{equation*}
\alpha_n(\psi_{n,\tau};\mathcal P)\leq
\varepsilon\end{equation*} 
and
\begin{equation*}
\beta_n(\psi_{n,\tau};\mathcal Q) \leq\exp\left\{(1-\lambda)\tau+n(\lambda-1)d\right\}.
\end{equation*}
\end{lemma}
Proof of lemma~\ref{lem:uniform-score-achievability} is given in Appendix~\ref{app:uniform-score-achievability}. Lemma~\ref{lem:uniform-score-achievability} reduces achievability to constructing a measurable function satisfying its two bounds. We now obtain such a function from a joint R\'enyi projection.

Let $\mu$ be a $\sigma$--finite measure dominating every law in $\mathcal P\cup\mathcal Q$, and define
\begin{equation*}
\mathcal D_\mu(\mathcal P) := \left\{\frac{dP}{d\mu}:P\in\mathcal P \right\},\qquad \mathcal D_\mu(\mathcal Q) := \left\{\frac{dQ}{d\mu}:Q\in\mathcal Q \right\}.
\end{equation*}
For $0<\lambda<1$, a pair $(Q_\lambda^\star,P_\lambda^\star)\in\mathcal Q\times\mathcal P$ is called a joint R\'enyi projection if
\begin{equation*}
Z_\lambda(Q_\lambda^\star,P_\lambda^\star) =\sup_{Q\in\mathcal Q} \sup_{P\in\mathcal P} Z_\lambda(Q,P).
\end{equation*}
We refer to $(Q_\lambda^\star,P_\lambda^\star)$ as the projected pair. Since $\lambda<1$, maximising $Z_\lambda(Q,P)$ is equivalent to minimising $D_\lambda(Q\|P)$. The following theorem gives conditions under which a projected pair exists and its log--likelihood ratio yields the required finite blocklength achievability bound.

%---------------------------------------------------
\begin{theorem}
\label{thm:projected-finite-blocklength-achievability} Fix $0<\lambda<1$. Suppose that $\mathcal D_\mu(\mathcal P)$ and $\mathcal D_\mu(\mathcal Q)$ are convex and weakly compact in $L^1(\mu)$. Then a projected pair $(Q_\lambda^\star,P_\lambda^\star)$ exists. Let $(Q_\lambda^\star,P_\lambda^\star)$ be a projected pair satisfying $Z_\lambda(Q_\lambda^\star,P_\lambda^\star)>0$ and $R\ll P_\lambda^\star+Q_\lambda^\star
\quad \forall R\in\mathcal P\cup\mathcal Q.$ Choose measurable versions
\begin{equation*}
a_\lambda^\star:=\frac{dP_\lambda^\star}{d(P_\lambda^\star+Q_\lambda^\star)}, \qquad b_\lambda^\star :=\frac{dQ_\lambda^\star}{d(P_\lambda^\star+Q_\lambda^\star)},
\end{equation*}
and define
\begin{equation*}
h_\lambda^\star(x):=\begin{cases}
\log\!\left(b_\lambda^\star(x)/a_\lambda^\star(x)
\right),
& a_\lambda^\star(x)>0,\ b_\lambda^\star(x)>0,\\
-\infty, & a_\lambda^\star(x)>0,\ b_\lambda^\star(x)=0, \\
+\infty, & a_\lambda^\star(x)=0,\ b_\lambda^\star(x)>0, \\
0, & a_\lambda^\star(x)=b_\lambda^\star(x)=0.
\end{cases}
\end{equation*}
Then
\begin{equation*}
D_\lambda(Q_\lambda^\star\|P_\lambda^\star) = D_\lambda(\mathcal Q\|\mathcal P),
\end{equation*}
and
\begin{equation*}
\begin{aligned}
\sup_{P\in\mathcal P} \mathbb E_P \left[e^{\lambda h_\lambda^\star} \right]&\leq e^{(\lambda-1)D_\lambda(\mathcal Q\|\mathcal P)}, \\ \sup_{Q\in\mathcal Q} \mathbb E_Q \left[ e^{(\lambda-1)h_\lambda^\star}\right] &\leq e^{(\lambda-1)D_\lambda(\mathcal Q\|\mathcal P)}.
\end{aligned}
\end{equation*}
Consequently, for every $n\geq1$ and $\varepsilon\in(0,1)$, the threshold test of Lemma~\ref{lem:uniform-score-achievability} with $h=h_\lambda^\star$, $d=D_\lambda(\mathcal Q\|\mathcal P)$, and
\begin{equation*}
\tau = \frac{ \log(1/\varepsilon) -n(1-\lambda)D_\lambda(\mathcal Q\|\mathcal P) }{\lambda
}
\end{equation*}
satisfies
\begin{equation*}
\alpha_n(\psi_{n,\tau};\mathcal P)\leq
\varepsilon
\end{equation*}
and
\begin{equation*}
\beta_n(\psi_{n,\tau};\mathcal Q)\leq \exp\left[-\frac{1-\lambda}{\lambda}\left(nD_\lambda(\mathcal Q\|\mathcal P)-\log\frac{1}{\varepsilon}
\right)\right].
\end{equation*}
Hence,
\begin{equation}
\begin{aligned}
\beta_n^\star(\varepsilon;\mathcal P,\mathcal Q) \leq \min\Bigg\{& 1-\varepsilon, \exp\left[ -\frac{1-\lambda}{\lambda} \left( nD_\lambda(\mathcal Q\|\mathcal P)-\log\frac{1}{\varepsilon}\right) \right]\Bigg\}.
\end{aligned}
\label{eq:projected-finite-blocklength-achievability}
\end{equation}

\end{theorem}


%##########################################
\appendices

\section{Proof of Theorem~\ref{thm:finite-blocklength-composite-renyi-converse}}
\label{app:finite-blocklength-composite-renyi-converse}
\begin{IEEEproof}
Fix a randomised test $\varphi_n$ satisfying
$\alpha_n(\varphi_n;\mathcal P)\leq\varepsilon$, and fix $P\in\mathcal P$, $Q\in\mathcal Q$, and $\lambda>1$. Suppose first that $D_\lambda(Q\|P)<+\infty$. Then $Q\ll P$, and hence $Q^{\otimes n}\ll P^{\otimes n}$. H\"older's inequality, together with $\varphi_n^{\lambda/(\lambda-1)}\leq\varphi_n$, the Type I constraint, and additivity of the R\'enyi divergence, gives
\begin{equation*}
\begin{aligned}
\mathbb E_{Q^{\otimes n}} &\left[\varphi_n(X^n) \right] = \mathbb E_{P^{\otimes n}}\left[\frac{dQ^{\otimes n}}{dP^{\otimes n}}(X^n) \varphi_n(X^n)\right]\\
&\leq\left(
\mathbb E_{P^{\otimes n}}
\left[\left(\frac{dQ^{\otimes n}}{dP^{\otimes n}}(X^n)\right)^\lambda\right]\right)^{1/\lambda}\left(\mathbb E_{P^{\otimes n}}
\left[\varphi_n(X^n)^{\lambda/(\lambda-1)}
\right]\right)^{(\lambda-1)/\lambda}\\
&\leq\exp\left\{\frac{n(\lambda-1)}{\lambda}
D_\lambda(Q\|P)\right\}\varepsilon^{(\lambda-1)/\lambda}=\exp\left\{-\frac{\lambda-1}{\lambda}
\left[\log\frac{1}{\varepsilon}-nD_\lambda(Q\|P)\right]\right\}.
\end{aligned}
\end{equation*}
Combining this inequality with
$\mathbb E_{Q^{\otimes n}}[\varphi_n(X^n)]\leq1$ yields
\begin{equation*}
\mathbb E_{Q^{\otimes n}}\left[1-\varphi_n(X^n)\right]\geq 1-\exp\left\{-\frac{\lambda-1}{\lambda}\left[\log\frac{1}{\varepsilon}-nD_\lambda(Q\|P)\right]_+\right\}.
\end{equation*}
If $D_\lambda(Q\|P)=+\infty$, the right-hand side is zero under the convention $[-\infty]_+=0$, so the same inequality follows from the nonnegativity of the Type II error probability. The preceding inequality holds for every $P\in\mathcal P$ and $Q\in\mathcal Q$. Its right-hand side is a continuous nonincreasing function of $D_\lambda(Q\|P)$. Therefore, the definition of $D_\lambda(\mathcal Q\|\mathcal P)$ gives
\begin{equation*}
\beta_n(\varphi_n;\mathcal Q)\geq 1- \exp\left\{-\frac{\lambda-1}{\lambda}\left[\log\frac{1}{\varepsilon}-nD_\lambda(\mathcal Q\|\mathcal P)\right]_+\right\},
\end{equation*}
without requiring the infimum over the pair to be attained. This bound holds for every $\lambda>1$. Since $x\mapsto1-e^{-x}$ is increasing, taking the supremum over $\lambda$ gives
\begin{equation*}
\beta_n(\varphi_n;\mathcal Q) \geq 1- \exp\left\{ -\sup_{\lambda>1} \frac{\lambda-1}{\lambda} \left[\log\frac{1}{\varepsilon} - nD_\lambda(\mathcal Q\|\mathcal P) \right]_+ \right\}.
\end{equation*}
Finally, taking the infimum over all randomised tests satisfying $\alpha_n(\varphi_n;\mathcal P)\leq\varepsilon$ proves
\eqref{eq:finite-blocklength-composite-renyi-converse}.
\end{IEEEproof}

%##################################################
\section{Proof of Lemma~\ref{lem:uniform-score-achievability}}
\label{app:uniform-score-achievability}

\begin{IEEEproof} The assumptions imply
\begin{equation*}
P\bigl(\{x\in\mathcal X:h(x)=+\infty\}\bigr)=0
\quad \forall P\in\mathcal P,
\qquad Q\bigl(\{x\in\mathcal X:h(x)=-\infty\}\bigr)=0
\quad \forall Q\in\mathcal Q.
\end{equation*}
Since $\lambda>0$, $e^{\lambda h(x)}=+\infty$ whenever $h(x)=+\infty$, whereas $\lambda-1<0$ implies $e^{(\lambda-1)h(x)}=+\infty$ whenever $h(x)=-\infty$. Thus, if either set had positive probability under the corresponding law, the relevant expectation in the assumptions would be infinite, contradicting the upper bound. The set of sequences on which both signs of infinity occur is measurable and is null under every $P^{\otimes n}$ and $Q^{\otimes n}$. Setting the sum to zero on this set therefore defines a measurable function $S_n:\mathcal X^n\to[-\infty,+\infty]$ without affecting either error probability. Since this convention changes the sum only on a null set,
the product structure gives
\begin{equation*}
\begin{aligned}
\sup_{P\in\mathcal P}\mathbb E_{P^{\otimes n}}\left[e^{\lambda S_n}\right]&=\sup_{P\in\mathcal P}\left(\mathbb E_P[e^{\lambda h}] \right)^n \leq e^{n(\lambda-1)d}, \\ \sup_{Q\in\mathcal Q} \mathbb E_{Q^{\otimes n}} \left[e^{(\lambda-1)S_n}\right] &= \sup_{Q\in\mathcal Q} \left( \mathbb E_Q[e^{(\lambda-1)h}] \right)^n \leq e^{n(\lambda-1)d}.
\end{aligned}
\end{equation*}
For any $\tau\in\mathbb R$, Markov's inequality gives
\begin{equation*}
\begin{aligned}
\alpha_n(\psi_{n,\tau};\mathcal P) &= \sup_{P\in\mathcal P} P^{\otimes n}\bigl(\{x^n\in\mathcal X^n:S_n(x^n)\geq\tau\}\bigr) \\ &\leq e^{-\lambda\tau+n(\lambda-1)d}, \\ \beta_n(\psi_{n,\tau};\mathcal Q) &= \sup_{Q\in\mathcal Q} Q^{\otimes n}\bigl(\{x^n\in\mathcal X^n:S_n(x^n)<\tau\}\bigr) \\ &= \sup_{Q\in\mathcal Q} Q^{\otimes n} \bigl( \{x^n\in\mathcal X^n: e^{(\lambda-1)S_n(x^n)}>e^{(\lambda-1)\tau}\} \bigr) \\ &\leq e^{(1-\lambda)\tau+n(\lambda-1)d}.
\end{aligned}
\end{equation*}

To satisfy the Type I constraint, it is sufficient that
\begin{equation*}
e^{-\lambda\tau+n(\lambda-1)d}\leq\varepsilon.
\end{equation*}
Since $\lambda>0$, this is equivalent to
\begin{equation*}
\tau \geq \frac{\log(1/\varepsilon)-n(1-\lambda)d}{\lambda},
\end{equation*}
which gives the required Type I and Type II bounds.
\end{IEEEproof}




%###################################################
\section{Proof of Theorem~\ref{thm:projected-finite-blocklength-achievability}}
\label{app:projected-finite-blocklength-achievability}
\begin{IEEEproof}
Let
\begin{equation*}
p_\lambda^\star
:=
\frac{dP_\lambda^\star}{d\mu},
\qquad
q_\lambda^\star
:=
\frac{dQ_\lambda^\star}{d\mu}.
\end{equation*}
For probability densities $p,\widetilde p,q,\widetilde q$ with respect
to $\mu$,
\begin{equation*}
\begin{aligned}
\left|
\int_{\mathcal X}
q^\lambda p^{1-\lambda}\,d\mu
-
\int_{\mathcal X}
\widetilde q^\lambda\widetilde p^{1-\lambda}\,d\mu
\right|
\leq
\|q-\widetilde q\|_1^\lambda
+
\|p-\widetilde p\|_1^{1-\lambda}.
\end{aligned}
\end{equation*}
Hence the Hellinger integral is continuous in the $L^1(\mu)$ norm.
Since the weighted geometric mean is jointly concave for
$0<\lambda<1$, its superlevel sets are convex and norm closed, and
therefore weakly closed. Thus, the Hellinger integral is weakly upper
semicontinuous on
$\mathcal D_\mu(\mathcal Q)\times\mathcal D_\mu(\mathcal P)$.
Weak compactness then guarantees the existence of a projected pair.

Since $\lambda-1<0$, maximising the Hellinger integral is equivalent to
minimising the R\'enyi divergence. Therefore,
\begin{equation*}
D_\lambda(Q_\lambda^\star\|P_\lambda^\star)
=
D_\lambda(\mathcal Q\|\mathcal P).
\end{equation*}

Fix $Q\in\mathcal Q$, and let $q=dQ/d\mu$. Convexity of
$\mathcal D_\mu(\mathcal Q)$ gives
\begin{equation*}
\int_{\mathcal X}
\left[
(1-t)q_\lambda^\star+tq
\right]^\lambda
(p_\lambda^\star)^{1-\lambda}\,d\mu
\leq
Z_\lambda(Q_\lambda^\star,P_\lambda^\star)
\end{equation*}
for every $0\leq t\leq1$. Applying
Lemma~\ref{lem:complete-feasible-directional-derivative} with
\begin{equation*}
\alpha=\lambda,
\qquad
w=(p_\lambda^\star)^{1-\lambda},
\qquad
y=q_\lambda^\star,
\qquad
x=q,
\end{equation*}
gives
\begin{equation*}
\mu\bigl(
\{x\in\mathcal X:
p_\lambda^\star(x)>0,\,
q_\lambda^\star(x)=0,\,
q(x)>0\}
\bigr)
=
0
\end{equation*}
and
\begin{equation*}
\begin{aligned}
\int_{\{x\in\mathcal X:
p_\lambda^\star(x)>0,\,
q_\lambda^\star(x)>0\}}
q(x)
\bigl(q_\lambda^\star(x)\bigr)^{\lambda-1}
\bigl(p_\lambda^\star(x)\bigr)^{1-\lambda}
\,d\mu(x)
\leq
Z_\lambda(Q_\lambda^\star,P_\lambda^\star).
\end{aligned}
\end{equation*}

Similarly, fix $P\in\mathcal P$, and let $p=dP/d\mu$. Convexity of
$\mathcal D_\mu(\mathcal P)$ gives
\begin{equation*}
\int_{\mathcal X}
(q_\lambda^\star)^\lambda
\left[
(1-t)p_\lambda^\star+tp
\right]^{1-\lambda}
\,d\mu
\leq
Z_\lambda(Q_\lambda^\star,P_\lambda^\star)
\end{equation*}
for every $0\leq t\leq1$. Applying
Lemma~\ref{lem:complete-feasible-directional-derivative} with
\begin{equation*}
\alpha=1-\lambda,
\qquad
w=(q_\lambda^\star)^\lambda,
\qquad
y=p_\lambda^\star,
\qquad
x=p,
\end{equation*}
gives
\begin{equation*}
\mu\bigl(
\{x\in\mathcal X:
q_\lambda^\star(x)>0,\,
p_\lambda^\star(x)=0,\,
p(x)>0\}
\bigr)
=
0
\end{equation*}
and
\begin{equation*}
\begin{aligned}
\int_{\{x\in\mathcal X:
p_\lambda^\star(x)>0,\,
q_\lambda^\star(x)>0\}}
p(x)
\bigl(q_\lambda^\star(x)\bigr)^\lambda
\bigl(p_\lambda^\star(x)\bigr)^{-\lambda}
\,d\mu(x)
\leq
Z_\lambda(Q_\lambda^\star,P_\lambda^\star).
\end{aligned}
\end{equation*}

The preceding support relations imply
\begin{equation*}
\begin{aligned}
P\bigl(
\{x\in\mathcal X:
p_\lambda^\star(x)=0,\,
q_\lambda^\star(x)>0\}
\bigr)
&=0
\quad \forall P\in\mathcal P,
\\
Q\bigl(
\{x\in\mathcal X:
p_\lambda^\star(x)>0,\,
q_\lambda^\star(x)=0\}
\bigr)
&=0
\quad \forall Q\in\mathcal Q.
\end{aligned}
\end{equation*}
Moreover, the domination assumption gives
\begin{equation*}
R\bigl(
\{x\in\mathcal X:
p_\lambda^\star(x)=0,\,
q_\lambda^\star(x)=0\}
\bigr)
=
0
\quad
\forall R\in\mathcal P\cup\mathcal Q.
\end{equation*}

The measurable versions in the theorem may be chosen so that
\begin{equation*}
a_\lambda^\star(x)
=
\frac{p_\lambda^\star(x)}
     {p_\lambda^\star(x)+q_\lambda^\star(x)},
\qquad
b_\lambda^\star(x)
=
\frac{q_\lambda^\star(x)}
     {p_\lambda^\star(x)+q_\lambda^\star(x)}
\end{equation*}
whenever
$p_\lambda^\star(x)+q_\lambda^\star(x)>0$, and both are zero otherwise.
Using the definition of $h_\lambda^\star$ and the preceding support
relations, for every $P\in\mathcal P$ and $Q\in\mathcal Q$,
\begin{equation*}
\begin{aligned}
\mathbb E_P
\left[
e^{\lambda h_\lambda^\star}
\right]
&=
\int_{\{x\in\mathcal X:
p_\lambda^\star(x)>0,\,
q_\lambda^\star(x)>0\}}
p(x)
\bigl(q_\lambda^\star(x)\bigr)^\lambda
\bigl(p_\lambda^\star(x)\bigr)^{-\lambda}
\,d\mu(x)
\\
&\leq
Z_\lambda(Q_\lambda^\star,P_\lambda^\star),
\\
\mathbb E_Q
\left[
e^{(\lambda-1)h_\lambda^\star}
\right]
&=
\int_{\{x\in\mathcal X:
p_\lambda^\star(x)>0,\,
q_\lambda^\star(x)>0\}}
q(x)
\bigl(q_\lambda^\star(x)\bigr)^{\lambda-1}
\bigl(p_\lambda^\star(x)\bigr)^{1-\lambda}
\,d\mu(x)
\\
&\leq
Z_\lambda(Q_\lambda^\star,P_\lambda^\star).
\end{aligned}
\end{equation*}
Taking the respective suprema and using
\begin{equation*}
Z_\lambda(Q_\lambda^\star,P_\lambda^\star)
=
e^{(\lambda-1)D_\lambda(\mathcal Q\|\mathcal P)}
\end{equation*}
gives the two bounds in the theorem.

Since
$Z_\lambda(Q_\lambda^\star,P_\lambda^\star)>0$,
the quantity $D_\lambda(\mathcal Q\|\mathcal P)$ is finite and
nonnegative. Lemma~\ref{lem:uniform-score-achievability}, applied with
\begin{equation*}
h=h_\lambda^\star,
\qquad
d=D_\lambda(\mathcal Q\|\mathcal P),
\end{equation*}
and
\begin{equation*}
\tau
=
\frac{
\log(1/\varepsilon)
-
n(1-\lambda)D_\lambda(\mathcal Q\|\mathcal P)
}{
\lambda
},
\end{equation*}
gives
\begin{equation*}
\alpha_n(\psi_{n,\tau};\mathcal P)
\leq
\varepsilon
\end{equation*}
and
\begin{equation*}
\beta_n(\psi_{n,\tau};\mathcal Q)
\leq
\exp\left[
-\frac{1-\lambda}{\lambda}
\left(
nD_\lambda(\mathcal Q\|\mathcal P)
-
\log\frac{1}{\varepsilon}
\right)
\right].
\end{equation*}
It follows that
\begin{equation*}
\beta_n^\star(\varepsilon;\mathcal P,\mathcal Q)
\leq
\exp\left[
-\frac{1-\lambda}{\lambda}
\left(
nD_\lambda(\mathcal Q\|\mathcal P)
-
\log\frac{1}{\varepsilon}
\right)
\right].
\end{equation*}

Finally, the constant randomised test
$\varphi_n(x^n)=\varepsilon$ for every $x^n\in\mathcal X^n$ has Type I
error $\varepsilon$ and Type II error $1-\varepsilon$. Taking the
smaller of the two bounds proves
\eqref{eq:projected-finite-blocklength-achievability}.

\end{IEEEproof}

\bibliographystyle{IEEEtran}
\bibliography{my}


\end{document}
